\documentclass[a4paper]{article}

\usepackage[a4paper, top=8mm, bottom=8mm, left=8mm, right=8mm]{geometry}

\usepackage{polyglossia}
\setdefaultlanguage[babelshorthands=true]{russian}
\setotherlanguage{english}

\usepackage{fontspec}
\setmainfont{FreeSerif}
\newfontfamily{\russianfonttt}[Scale=0.7]{DejaVuSansMono}

\usepackage[tiny, compact]{titlesec}

\usepackage{titling}
\setlength{\droptitle}{-1cm}
\pretitle{\begin{center}\begin{bfseries}\Large}
\posttitle{\par\end{bfseries}\end{center}}
\preauthor{\begin{center}\normalsize}
\postauthor{\par\end{center}\vspace{-1.8cm}}

\usepackage{hyperref}
\usepackage{bookmark}
\usepackage{csquotes}
% Сюда можно дописывать нужные вам пакеты.

\title{Оптимизация алгоритмов дельта-кодирования GDelta в системе MarLine}

\author{Исаев Даниил Сергеевич}

% Дата много места занимает, её лучше не указывать.
\date{}

\begin{document}

\maketitle

\begin{flushright}
    Группа: \emph{25.Б72-мм}

    Кафедра: \emph{Кафедра системного программирования СПбГУ}

    % Степень/звание и должность научника мы лучше вас знаем, их можно не указывать.
    Научный руководитель: \emph{Гориховский Вячеслав Игоревич}

    % А вот про консультанта укажите официальное место работы, с "ООО" или чем-то таким, и должностью желательно по трудовой, не просто "techlead". Выясните у консультанта.
    Консультант: \emph{Дмитриевцев Алексей Сергеевич}
    
    Младший инженер по разработке ПО ООО "КНС ГРУПП"
    
    Номер семестра практики: \emph{3}
\end{flushright} 


\section{Постановка задачи}

Работа выполняется в рамках open source-проекта MarLine --- системы хранения данных с двухуровневой дедупликацией: контентно-ориентированное разбиение (CDC) отсекает неизменные данные, а Similarity-Based Chunking подбирает для каждого чанка похожий базовый чанк и хранит его в виде дельты, полученной алгоритмом GDelta. \textbf{Целью работы является} оптимизация подсистемы дельта-кодирования GDelta в системе MarLine: сократить время выбора базового чанка, время формирования скетча чанка и время его дельта-кодирования.

Для достижения цели требуется решить следующие задачи:

\begin{enumerate}
    \item Спроектировать модель представления чанка, сочетающую 64-битный скетч и 32-битный признак для поиска базового чанка, совместимую с rolling-hash, применяемым в GDelta.
    \item Реализовать модуль генерирования скетча чанка и включить его в общий конвейер обработки.
    \item Реализовать префильтр выбора базового чанка и устранить линейный перебор кандидатов.
    \item Внедрить в алгоритм GDelta правило исключения: пропуск обращения к Fp-таблице при различии битов класса слова в скетчах входного и базового чанков.
\end{enumerate}

Результаты работы необходимы для снижения накладных расходов подсистемы сжатия: чем быстрее кодируются чанки, тем меньше времени и процессорных ресурсов уходит на запись и разбор данных. В выполнении работы заинтересованы команда разработки проекта MarLine.

\section{Обзор}

\section{Описание предлагаемого решения}

% Раздел пишется в свободной форме, с минимумом технических подробностей (можно приводить ссылки на более подробные документы).

%Тут надо кратко описать идею решения, спозиционировать его относительно предыдущих наработок\footnote{Если на них надо сослаться, то в подстраничной сноске, URL: \url{http://example.com} (дата обращения: не забываем указывать).} и существующих аналогов (тоже предельно кратко), по возможности обосновать принятые решения.

% Вот это важно, опишите кратко, что использовали
%Реализация решения осуществлялась с использованием \dots (тут описать используемые технологии).

\section{Апробация / Эксперименты}

% Название секции можно поменять на более подходящее вашей работе

%Название этой секции --- либо \enquote{Апробация}, либо \enquote{эксперименты} (если содержательно есть и то и то, можно две отдельные секции сделать). Тут опишите, как проводилась проверка того, что получилось что-то адекватное. Если есть отзывы пользователей, отлично (особенно если они в структурированном виде, типа анкеты SUS). Если есть отзывы команды, консультанта и т.п., хорошо. Если есть только тесты. то хотя бы их тут опишите.

%Если работа подразумевала замеры чего-то (например, производительности), приведите тут итоговые результаты и выводы, и дайте ссылку на таблицы с данными. Не забывайте про матстат (приводите матожидание и среднеквадратичное отклонение, не просто результат замера).

%Приведите выводы из эксперимента.

\section{Заключение}

%В ходе выполнения работы получены следующие результаты:

% ниже списком перечисляете то, что выносится как результат практики.

%\begin{itemize}
%    \item \dots
%    \item \dots
%    \item \dots
%\end{itemize}

%Материалы, иллюстрирующие выполнение работы, размещены по адресу: \url{http://www.example.com} (если есть и применимо --- скриншоты, видео на 1-2 минуты и т.п. с демонстрацией; может быть очень полезно даже для консольных утилит или библиотек).

%Репозиторий/ссылка на пуллреквест: \url{http://www.example.com} (это обязательно; ссылок может быть несколько).

%Техническая документация: \url{http://www.example.com} (если применимо --- для пуллреквестов можно техническую сторону прямо в комментарии к пуллреквесту указать, но можно и в отдельную вики-страницу вынести).

%Решение внедрено в производственные процессы компании/прошло процесс ревью/находится на ревью/находится в разработке и т.п. В общем, кратко сформулируйте, что в итоге получилось и в каком статусе результат, чтобы не было ощущения, что оно сделано, положено на GitHub и забыто.

\end{document}
